%%%%%%%%%%%%%%%%%%%%%%%%%%%%%%%%%%%%%%%%%%%%%%%%%%%%%%%%%%%%
% Filename: Pu_Theory_Moebius_Ring_Neutrino_Chirality_Functor_v11_en.tex
% Compile: xelatex or pdflatex
% Encoding: UTF-8
% Positioning: A formalized statement within the Pu-theory framework —
%              a functor derivation, not a standard physics paper
% v11 (EN): Translation of v11. Axioms I & II declared at Definition 1.1;
%           conclusion item 2 anchored; all else preserved verbatim.
%%%%%%%%%%%%%%%%%%%%%%%%%%%%%%%%%%%%%%%%%%%%%%%%%%%%%%%%%%%%
\documentclass[11pt]{article}

\usepackage{amsmath,amssymb,amsthm,mathrsfs}
\usepackage[margin=2.5cm]{geometry}
\usepackage[utf8]{inputenc}
\usepackage[T1]{fontenc}
\usepackage{hyperref}

\theoremstyle{plain}
\newtheorem{thm}{Theorem}[section]
\newtheorem{prop}[thm]{Proposition}
\newtheorem{cor}[thm]{Corollary}

\theoremstyle{definition}
\newtheorem{defn}[thm]{Definition}
\newtheorem{rem}[thm]{Remark}

\newcommand{\Cat}{\mathbf}

\begin{document}

\title{Functorial Correspondence between M\"obius-Ring One-Sidedness
and Neutrino Left-Handedness\\
{\large --- Pu-Theory Formalization Interface, v11: Complex-Exponent Extension}}
\author{Author}
\date{October 14, 2026}
\maketitle

\begin{abstract}
Building on the chirality-functor derivation of v8, this paper upgrades the
source-category point-seed from the real exponent $0^{p/q}$ to the complex
exponent $0^{z}$ ($z\in\mathbb{C}$). Core results:
\textbf{Real case} ($z=p/q\in\mathbb{R}$) inherits v8 verbatim --- neutrino
left-handedness is the necessary image of the functor $F^{p/q}$, $q=2$
returns the M\"obius ring, $q=3$ yields the three-generation flavor
structure; \textbf{Complex case} ($z=a+ib$, $b\neq 0$) yields, within the
same functor framework, structural interfaces for the oscillation phase and
CP violation. The paper explicitly separates three tiers: proved (real
case), structurally fixed (complex case), and open interfaces (quantitative
details), preserving entry points for later extension of the Pu-theory
framework.
\end{abstract}

\section{Source Category: Complex-Exponent Point-Seed}

\begin{defn}[Point-seed $0^z$]\label{def:pointseed}
Two primitive axioms of Pu-theory provide the foundation of this definition:
\textbf{Axiom I}: inside the point-seed, $0=1$ (internal identity, source of
asymmetry); \textbf{Axiom II}: $1/0=\infty$ (foundation of the formal
extension).

Below, $0^z := e^{z\ln 0}$ is a \textbf{formal extension within the Pu-theory
framework}; convergence in classical complex analysis is not required;
$\ln 0 = -\infty$ is supplied by Axiom II and used as a formal symbol, not
treated as a theorem of classical complex analysis. For
$z = a+ib$ ($a,b\in\mathbb{R}$), the classical extension factors as:
\[
0^{a+ib} = 0^a\cdot 0^{ib},\qquad 0^{ib} = e^{-ib\infty}.
\]
For $b\neq 0$, $0^{ib}$ is a formal phase symbol, not an object of classical
complex analysis; its physical scale is given by the extension interface in
\S\ref{sec:extension}.
\end{defn}

\begin{rem}[Status of real and imaginary parts]\label{rem:ab}
$a$ and $b$ are the real and imaginary parts of the same complex number
$z$; they are on an equal footing, with no primary/secondary distinction.
In the functor framework:
\begin{itemize}
\item Real part $a$: yields the $q$-branch structure (when $a = p/q$ is
rational), mapped by the functor to the flavor structure;
\item Imaginary part $b$: yields the phase factor $e^{-ib\infty}$, mapped by
the functor to the oscillation phase.
\end{itemize}
This division is not an additional convention but the two readings of the
complex structure under the functor.
\end{rem}

\begin{defn}[Source category $\Cat{Pu}_z$]\label{def:puz}
Objects: geometric structures generated by the complex exponent $0^z$.
\begin{itemize}
\item $D^z$: untwisted complex band ($z$ determines the branch structure
and phase structure; locally two-sided / globally two-sided)
\item $M^z$: twisted complex band (as above; locally two-sided / globally
one-sided)
\end{itemize}
Generating morphism:
\[
t^z: D^z \to M^z
\]
(executes one twist on each branch, preserving the $\mathbb{Z}_q$ cyclic
symmetry and the phase structure). Identity morphisms and composition rules
agree with $\Cat{Pu}$ of v8. $\Cat{Pu}_z$ is the free category generated by
the objects $\{D^z,M^z\}$ and the generating morphism $\{t^z\}$.
\end{defn}

\begin{rem}[$z\in\mathbb{R}$ returns to v8]\label{rem:backtov8}
For $z = p/q\in\mathbb{R}$ ($b=0$):
\begin{itemize}
\item $q=2$: the uniformization of $0^{1/2}$ is the M\"obius ring,
$\Cat{Pu}_{1/2} = \Cat{Pu}$ (the source category of v8);
\item $q=3$: $0^{1/3}$ yields a three-branch structure, corresponding to
three generations of flavor;
\item general $q$: $0^{1/q}$ yields $q$ branches.
\end{itemize}
\end{rem}

\begin{rem}[$z\in\mathbb{C}$ adds a phase structure]\label{rem:complexcase}
For $z = a+ib$ ($b\neq 0$):
\begin{itemize}
\item $a = p/q$ yields the $q$-branch structure (as in the real case);
\item $b$ yields the phase $e^{-ib\infty}$, appearing in $\Cat{Pu}_z$ as a
$\mathrm{U}(1)$ connection on the complex band.
\end{itemize}
\end{rem}

\section{Target Category: Fermions and Weak Interaction}

\begin{defn}[Category $\Cat{Ferm}$]\label{def:ferm}
Objects: fermion spinor structures.
\begin{itemize}
\item $\Psi$: the initial state under the action of the weak-interaction
vertex, a fermion field with left- and right-handed chiral components
(two-state; locally two-state / globally two-state)
\item $\Psi_L$: the fermion field after weak interaction (only the
left-handed component survives; one-state; locally two-state / globally
one-state)
\end{itemize}
Morphisms: chiral operations.
\begin{itemize}
\item $w: \Psi \to \Psi_L$: weak-interaction vertex (projects a two-state
fermion onto a one-state chiral fermion)
\end{itemize}
Identity morphisms and composition rules agree with v8. Projection
requirements:
\begin{enumerate}
\item (Idempotence) there exists an inclusion morphism
$\iota: \Psi_L \hookrightarrow \Psi$ such that
$w\circ\iota = \mathrm{id}_{\Psi_L}$.
\item (Completeness) the fixed-point set $\mathrm{Fix}(w) = \Psi_L$,
i.e.\ $\Psi\setminus\Psi_L$ contains no fixed points.
\end{enumerate}
\end{defn}

\begin{rem}[Physical correspondence]\label{rem:physical}
$\Psi$ corresponds to the initial state, under the action of the
weak-interaction vertex, of fermion fields participating in $SU(2)_L$ weak
interaction (e.g.\ the description in which the lepton doublet
$(\nu_L,e_L)$ coexists with the right-handed companion field $e_R$).
$\Psi_L$ corresponds to the left-handed component that alone couples to the
weak interaction. $w$ corresponds to the chirality projection at the
weak-interaction vertex.
\end{rem}

\begin{rem}[Target uniqueness is definitional, not a theorem]\label{rem:targetbydef}
In $\Cat{Ferm}$, the only non-identity morphism out of $\Psi$ is, by
definition, $w$; hence the target $\Psi_L$ of $w$ is given directly by the
definition. That the weak-interaction vertex is unique (the $SU(2)_L$
chirality projection) is a physical input of the minimal Standard Model,
not a mathematical consequence of $\Cat{Ferm}$.
\end{rem}

\begin{rem}[Fermions that have not traversed the twist]\label{rem:untwisted}
The right-handed electron $e_R$ does not participate in $SU(2)_L$ weak
interaction: it carries only $U(1)_Y$ hypercharge ($Y=-2$; we adopt the
convention $Q = T_3 + Y/2$), acquiring a partial coupling to the $Z$ boson
via electroweak mixing, but it does not enter the $SU(2)_L$ doublet. This
corresponds to the complex band $D^z$ in $\Cat{Pu}_z$ that has not undergone
the twist. If a sterile $\nu_R$ is introduced, the same reasoning
applies --- it does not participate in weak interaction and does not
undergo the chirality projection.
\end{rem}

\section{The Functor}

\begin{defn}[Complex functor $F^z: \Cat{Pu}_z \to \Cat{Ferm}$]\label{def:functor}
Object map:
\begin{align*}
F^z(D^z) &= \Psi \\
F^z(M^z) &= \Psi_L
\end{align*}
Morphism map:
\begin{align*}
F^z(t^z) &= w \\
F^z(\mathrm{id}_{D^z}) &= \mathrm{id}_\Psi \\
F^z(\mathrm{id}_{M^z}) &= \mathrm{id}_{\Psi_L}
\end{align*}
Composition is preserved as in v8:
$F^z(t^z\circ\mathrm{id}_{D^z}) =
F^z(t^z)\circ F^z(\mathrm{id}_{D^z}) = w$.
\end{defn}

\begin{rem}[The functor preserves structure]\label{rem:preserve}
$F^z$ maps the source $D^z$ of $t^z$ (two-sided structure) to $\Psi$
(two-state structure), and the target $M^z$ (one-sided structure) to
$\Psi_L$ (one-state structure). For $z\in\mathbb{R}$, $F^z$ agrees
completely with $F$ of v8. For $z\in\mathbb{C}$, $F^z$ additionally maps
the $\mathrm{U}(1)$ structure of the complex band to the phase structure
on $\Psi_L$ (see \S\ref{sec:osc}).
\end{rem}

\begin{rem}[$z=p/q$ inherits v8 verbatim]\label{rem:inherit}
For $z = p/q\in\mathbb{R}$, $F^{p/q} = F$ (the functor of v8); $q=2$
returns the M\"obius ring of v8, and $q=3$ yields the three-generation
flavor structure.
\end{rem}

\section{Fixing of the Left-Handed Direction}

\begin{defn}[Correspondence between point-seed components and $SU(2)_L$ weak isospin]\label{def:su2}
The two components of the point-seed (the real case of $0^{p/q}$ at $q=2$)
correspond to the third component $T_3$ of weak isospin in the $SU(2)_L$
doublet as follows:

The first component corresponds to the state with $T_3 = +1/2$
($\nu_L$; from the charge formula $Q = T_3 + Y/2$ with $Y=-1$, giving
$Q=0$);
the second component corresponds to the state with $T_3 = -1/2$ ($e_L$,
giving $Q=-1$).

In the lepton doublet basis
$\begin{pmatrix}\nu_L\\ e_L\end{pmatrix}$, the upper component $\nu_L$
corresponds to $T_3=+1/2$ and the lower component $e_L$ corresponds to
$T_3=-1/2$.

(Note: that $\nu_L$ occupies the upper component of the doublet with
$T_3=+1/2$ is a convention of the Standard Model; under the opposite
ordering, the correspondence in this definition must be flipped
accordingly.)
\end{defn}

\begin{rem}[Independence of chirality and isospin]\label{rem:independence}
Definition~\ref{def:su2} establishes a correspondence between point-seed
components and $T_3$, but this is a \emph{correspondence convention} between
the Pu-theory point-seed and the $SU(2)_L$ doublet, not a conclusion
derived from Pu-theory itself. Pu-theory only supplies the fact that the
point-seed has two components; the choice that the first component
corresponds to $T_3=+1/2$ rather than $-1/2$ is a choice of binary
correspondence $0\leftrightarrow +1/2$, $1\leftrightarrow -1/2$ --- under
the opposite choice, Pu-theory remains internally consistent.

Moreover, $T_3$ and chirality ($P_L/P_R$) are two independent quantum
numbers. The representation structure of $SU(2)_L$ does not by itself
distinguish left from right chirality --- it specifies only the $T_3$ values
of the doublet. The fact that ``the components of the lepton doublet are
left-handed'' is an experimental input from the minimal Standard Model.

Hence the mechanism by which the functor $F^z$ maps the directionality of
the twist to $P_L$ is \emph{a joint effect of two tiers}:
\begin{enumerate}
\item The correspondence-convention tier (Pu-theory side): the
correspondence between point-seed components and $SU(2)_L$ doublet $T_3$;
\item The empirical tier (physics side): in the minimal Standard Model,
weak interaction couples only to left-handed fermions.
\end{enumerate}
The two tiers combine and fix the functor direction.
\end{rem}

\section{Chirality Theorem (Real Case)}

\begin{thm}[Functor necessity given the minimal Standard Model structure]\label{thm:main}
Let $z = p/q\in\mathbb{R}$. In $\Cat{Pu}_{p/q}$, the asymmetry of the
point-seed $0^{p/q}$ generates the twist morphism
$t^{p/q}: D^{p/q}\to M^{p/q}$, which is the unique generating morphism from
``two sides'' to ``one side'' in the freely generated category.

By the functor $F^{p/q}$ (Definition~\ref{def:functor}),
$F^{p/q}(t^{p/q}) = w$ is the corresponding ``two-state $\to$ one-state''
morphism in $\Cat{Ferm}$, whose target is given by definition as $\Psi_L$.

By Definition~\ref{def:su2} and Remark~\ref{rem:independence}, the
Pu-theory--$SU(2)_L$ correspondence convention together with the empirical
facts of the minimal Standard Model jointly fix the direction as the
left-handed projection $P_L$.

Therefore, given the functor $F^{p/q}$ and the minimal Standard Model
structure, neutrino left-handedness is the necessary result of the
M\"obius-ring twist structure being transmitted by the functor to the
fermion category.
\end{thm}

\begin{proof}
Verbatim identical to the proof of Theorem 6.1 of v8:
the asymmetry of the point-seed $0^{p/q}$ generates the unique morphism
$t^{p/q}: D^{p/q}\to M^{p/q}$ in the free category $\Cat{Pu}_{p/q}$;
the functor $F^{p/q}$ preserves this structure, $F^{p/q}(t^{p/q}) = w$ maps
the two-state object $\Psi$ to the one-state object $\Psi_L$, the target
being given by definition as noted in Remark~\ref{rem:targetbydef};
directionality is fixed by the two tiers (Definition~\ref{def:su2} and
Remark~\ref{rem:independence}), determining $w = P_L$;
in the minimal Standard Model, the neutrino enters the lepton doublet
$(\nu_L,e_L)$ only through its left-handed component, so the instance of
$\Psi_L$ in the lepton sector is precisely $\nu_L$.
\end{proof}

\begin{rem}[$q=2$ returns v8; $q=3$ gives three generations]\label{rem:qcases}
For $p/q = 1/2$, Theorem~\ref{thm:main} is the neutrino left-handedness
theorem of v8. For $p/q = 1/3$, $\Psi_L$ automatically acquires a
$\mathbb{Z}_3$ flavor structure; the three generations of flavor come from
the three branches of $0^{1/3}$; see \S\ref{sec:flavor}.
\end{rem}

\section{Flavor Structure: Real-Part Branches}\label{sec:flavor}

\begin{rem}[Flavor comes from the branch structure of $a$]\label{rem:flavor}
The real part $a = p/q$ of $z$ yields the $q$-branch structure, mapped by
$F^{p/q}$ to the $\mathbb{Z}_q$ flavor structure of $\Psi_L$:
\[
\Psi_L = \bigoplus_{k=0}^{q-1} \Psi_L^{(k)}.
\]
Flavor eigenstates are not fixed points of the branch permutation but
eigenstates of the $\mathbb{Z}_q$ branch permutation. $q=3$ yields three
generations of flavor.
\end{rem}

\begin{rem}[The flavor structure is an intrinsic property of the functor image]\label{rem:flavorfrom}
The $\mathbb{Z}_q$ flavor structure of $\Psi_L$ is not an additional input;
it is the image transmitted by $F^{p/q}$ directly from the $q$-branch
structure of $M^{p/q}$. Each branch of $\Cat{Pu}_{p/q}$ corresponds to each
flavor component of $\Psi_L$.
\end{rem}

\section{Oscillation and CP: Imaginary-Part Phase}\label{sec:osc}

\begin{rem}[Oscillation: the phase structure of $b$]\label{rem:osc}
The imaginary part $b$ of $z = a+ib$ is mapped by $F^z$ to the phase factor
\[
e^{-ib\infty},
\]
which in $\Cat{Ferm}$ manifests as the complex phase of $\Psi_L$.
In the flavor basis, the oscillation morphism takes the form
\[
\mathrm{osc}: \Psi_L^{(k)} \longmapsto \sum_{l=0}^{q-1}
U_{kl}\, e^{i\theta_l}\,\Psi_L^{(l)},
\]
where the \textbf{structural origin} of $U_{kl}$ is the $\mathbb{Z}_q$
branch permutation, and $e^{i\theta_l}$ is the phase accumulation of branch
$l$ (the image of $b$).

\textbf{Embedding remark.} The regular representation of $\mathbb{Z}_q$ is
the $q\times q$ cyclic permutation matrix, a discrete subgroup of
$\mathrm{U}(q)$ that does not contain continuous mixing angles. The precise
embedding from the $\mathbb{Z}_q$ branch permutation to the physically
relevant $SU(q)$ mixing matrix ($q=3$: the $3\times 3$ unitary PMNS matrix)
\textbf{is left to later extension} (see \S\ref{sec:extension}); the
present version supplies only that the structural origin of $U_{kl}$ is the
branch permutation, and does not claim to have produced the continuous
parameters of $SU(q)$.

The quantitative correspondence between PMNS mixing angles and oscillation
frequencies is also left to later extension.
\end{rem}

\begin{rem}[CP violation: non-reality of the complex exponent]\label{rem:cp}
The CP transformation corresponds on the source side to $z\mapsto\bar z$
(complex conjugation). For $b\neq 0$, $0^z\neq 0^{\bar z}$, i.e.\ the
complex exponent is non-real.

\textbf{Functor bridge.} On the target side, the CP transformation
corresponds to complex conjugation of the PMNS matrix,
$U_{PMNS}\mapsto U_{PMNS}^*$. For $b\neq 0$, $F^z\neq F^{\bar z}$; hence,
after transmission through the functor, $U_{PMNS}$ cannot be absorbed into
a real orthogonal matrix --- its complex phase is irremovable, which is CP
violation.

For $q\geq 3$, the complex characters $\omega^k$ and $\omega^{-k}$ of
$\mathbb{Z}_q$ are unequal, consistent with the irremovability of the
CP-violating phase. For $q=2$, $\mathbb{Z}_2$ is self-conjugate,
$U_{PMNS}$ can be absorbed into a real orthogonal matrix, and there is no
CP violation. Quantitative details (the exact values of the PMNS phase
angles) are left to later extension.
\end{rem}

\begin{rem}[$q\geq 3$ and physical consistency with CP violation]\label{rem:km}
Physically, CP violation in the Standard Model requires at least three
generations (Kobayashi--Maskawa). In the present framework, $q\geq 3$
renders the $\mathbb{Z}_q$ branch structure non-self-conjugate, consistent
with the physical condition.
\end{rem}

\section{Extension Interface}\label{sec:extension}

\begin{rem}[Three-tier boundary of the present version]\label{rem:boundary}
This version (v11) is complete in the following sense:
\begin{enumerate}
\item \textbf{Proved} ($z\in\mathbb{R}$): all theorems of the real case
hold verbatim (everything up to \S\ref{sec:flavor}, and
Theorem~\ref{thm:main}).
\item \textbf{Structurally fixed} ($z\in\mathbb{C}$, $b\neq 0$):
$0^z = 0^a\cdot 0^{ib}$ yields the two-tier real/imaginary structure; the
real part yields branches (flavor), the imaginary part yields phase
(oscillation), and the functor $F^z$ maps these respectively to the flavor
structure and phase structure of $\Psi_L$.
\item \textbf{Open interfaces} (quantitative details):
  \begin{itemize}
  \item \textbf{Embedding of $\mathbb{Z}_q$ into $SU(q)$}:
  the precise manner of embedding from the branch permutation into the
  continuous PMNS mixing matrix (for $q=3$: three mixing angles and one CP
  phase);
  \item Physical scale of $b$: the quantitative correspondence between the
  oscillation frequency $\theta_l$ and the mass-squared difference
  $\Delta m^2$, baseline $L$, and energy $E$;
  \item Rational-approximation rule for $a$: how the branch number $q$
  emerges when $a$ is continuous;
  \item Periodic structure of $\infty$ under the exponential map: how the
  $2\pi$-periodicity of $e^{-ib\infty}$ is derived precisely from
  Axiom III.
  \end{itemize}
\end{enumerate}
These are not deficiencies of v11 but the design of Pu-theory as an
extensible framework: the present version supplies the structure and the
proved part; the details are completed by later extension in line with the
physical interface.
\end{rem}

\begin{rem}[Methodological note]\label{rem:method}
This version continues the v8 approach: explicitly mark what is proved,
what is convention, and what is left to later extension. v8 marked ``the
$T_3$ correspondence is a convention''; v11 additionally marks ``the scale
of $b$ is an interface'' and ``the embedding $\mathbb{Z}_q\to SU(q)$ is an
interface''. This posture is consistent with the Pu-theory framework
``nature has no forbidden zones'': the framework is not rejected because of
open interfaces, nor is anything unproved treated as proved.
\end{rem}

\section{Conclusion}

\begin{enumerate}
\item Three generations of flavor come from the real branches of $0^{1/3}$
($a = 1/3$).
\item Neutrino left-handedness comes from the internal identity
(Axiom I: $0=1$ inside the point-seed) of each branch of $0^{p/q}$
(real case, Theorem~\ref{thm:main}).
\item Oscillation and CP come from the complex structure of $0^z$
($b\neq 0$), mapped by the functor $F^z$ to the phase structure of
$\Psi_L$; quantitative details are left to later extension.
\item All structure comes from the classical extension of
$0^z = e^{z\ln 0}$.
\item The real case is closed (all theorems of v8 are inherited verbatim);
the complex case has open interfaces (\S\ref{sec:extension}) --- in keeping
with the Pu-theory framework ``nature has no forbidden zones''.
\end{enumerate}

\appendix
\section{Appendix: Translation between Morphisms and Objects}

In category theory, a morphism $f: A\to B$ can be turned into an object in
the following ways:
\begin{enumerate}
\item \textbf{Arrow category} $\mathrm{Arr}(\mathcal{C})$: its objects are
morphisms $f: A\to B$ of $\mathcal{C}$, its morphisms are commuting
squares. In $\mathrm{Arr}(\Cat{Pu}_z)$, $t^z$ is itself an object.
\item \textbf{Image object}: $\mathrm{im}(t^z) = M^z$ as an object of
$\Cat{Pu}_z$.
\item \textbf{Slice category} $\mathcal{C}/M^z$: objects are all morphisms
with target $M^z$.
\end{enumerate}

The functor $F^z$ is defined as an ordinary functor on $\Cat{Pu}_z$; $F^z$
induces a functor
$\mathrm{Arr}(F^z): \mathrm{Arr}(\Cat{Pu}_z) \to \mathrm{Arr}(\Cat{Ferm})$,
mapping the object $t^z$ to the object $w$.

\section{Appendix: Correspondence Table from v8 to v11}

\begin{center}
\begin{tabular}{lll}
\hline
\textbf{Structure} & \textbf{v8} & \textbf{v11} \\
\hline
Point-seed & $\vartheta = (0,1)$ & $0^z$, $z\in\mathbb{C}$ \\
Source category & $\Cat{Pu}$ ($q=2$) & $\Cat{Pu}_z$ ($q$ determined by $a$) \\
Generating morphism & $t: D\to M$ & $t^z: D^z\to M^z$ \\
Target category & $\Cat{Ferm}$ & $\Cat{Ferm}$ (unchanged) \\
Functor & $F$ & $F^z$ \\
Flavor structure & (none) & $\Psi_L = \bigoplus_{k=0}^{q-1}\Psi_L^{(k)}$ \\
Phase structure & (none) & $e^{-ib\infty}$ \\
Chirality theorem & Theorem 6.1 & Theorem~\ref{thm:main} (real case) \\
\hline
\end{tabular}
\end{center}

For $z = 1/2$ ($b=0$), v11 agrees with v8 completely.

\end{document}